\documentclass[10pt,a4paper]{amsart}
\usepackage[margin=19mm]{geometry}
\usepackage{amsmath,amssymb,mathtools}
\usepackage[scr=rsfs,cal=euler]{mathalfa}
\usepackage{xfrac}
\usepackage[unicode,hidelinks,bookmarks=false]{hyperref}
\newcommand{\ie}{i.e.\ }
\newcommand{\cf}{cf.\ }
\newcommand{\resp}{respectively\ }
\newcommand{\Subsection}{\subsection}
\providecommand{\nobreakdash}{\nobreak\hbox{-}}
\setlength{\emergencystretch}{2em}
\multlinegap0pt
\usepackage{bm}
\usepackage{bbm}
\usepackage{dsfont}
\usepackage{xspace}
\usepackage{cancel}
\usepackage{mathtools}
\usepackage{amsthm}
\makeatletter
\@ifundefined{theorem}{\newtheorem{theorem}{Theorem}}{}
\@ifundefined{lemma}{\newtheorem{lemma}[theorem]{Lemma}}{}
\makeatother
\title[Martingale Problem and Quadratic Family]{Martingale Problem and Quadratic Family}
\author{Haoming Wang}
\date{Source: arXiv:2509.06016v1 (CC BY 4.0)}
\begin{document}
\maketitle

\noindent\emph{Source and licence.} Martingale Problem and Quadratic Family. Version arXiv:2509.06016v1 (CC BY 4.0), distributed under the Creative Commons Attribution 4.0 International licence, as recorded in the arXiv licence metadata for this exact version and shown on the abstract page https://arxiv.org/abs/2509.06016v1. Full rights notice: licenses/arxiv-2509.06016-CC-BY-4.0.txt.\par\smallskip

\noindent\textit{Editorial scope.} Counted unit: the Theorem whose source label is \texttt{None} in the article (source0.tex lines 102--108, source label \texttt{None}), delivered together with the article's own complete original proof of it (source0.tex lines 109--196, one \texttt{proof} environment of 88 lines). No mathematics is written, repaired or re-derived: statement and proof are the article's own text line for line. Required context retained uncounted: none. Every definition, display and label of those lines is retained unchanged. Omitted from this package and not redistributed: 1--101, 197--337. Nothing inside a retained excerpt is omitted or altered: every retained body is delivered exactly as the article's lines are written, so no omission receipt of any kind accompanies this package. Citations: the retained excerpts carry no citation command at all, so no bibliography entry of the article is transcribed and no reference is redistributed. Reference adaptation: none was needed and none was made; every reference inside the delivered unit resolves inside it, so no unresolved reference or citation is produced. Printed slips kept literal: no printed word, symbol, hypothesis or number of the counted statement or of its proof was altered. Layout and class: the article's own class is \texttt{lms} with packages including amsfonts, amsmath, amssymb, mathrsfs, natbib; the wrapper is the lane's pinned \texttt{amsart} editorial wrapper at 10pt on A4, which restates the theorem-like environments the retained text needs and adds the article's own macro definitions that the retained text needs (none), with extra wrapper packages (bm, bbm, dsfont, xspace, cancel, mathtools). The delivered numbering is the unit's own. Assets: no retained span contains a figure environment, a graphics-inclusion command, a PDF-page inclusion, a table or a TikZ picture; no image file is copied into this package, the package declares no asset dependency and no image file is redistributed with it. Licence: version arXiv:2509.06016v1 of this article is distributed under the Creative Commons Attribution 4.0 International licence, as recorded in the arXiv licence metadata for this exact version and shown on the abstract page https://arxiv.org/abs/2509.06016v1. A full rights notice is delivered at licenses/arxiv-2509.06016-CC-BY-4.0.txt. Characters: every retained excerpt is pure ASCII, so the delivered engine, XeLaTeX with its default Latin Modern text font, reports no missing character.
\begin{theorem}
    Let $X_{n}$ be a Markov chain with $\sigma(X_{0},X_{1},\dots, X_{n})$-predictable transition probability matrix $P_n$ under $\mathbf P$. Suppose $\mathbf{P}_{0}$ is a probability measure with respect to which $X_{n}$ is the stationary Markov chain with transition probability matrix $P_0$ satisfying $P_{0}(i.j) >0$ and the same initial distribution, then $\mathbf{P} \ll \mathbf{P}_{0}$ and the Girsanove transformation (D) holds,
    \begin{equation*}   \left.d\mathbf{P}/d\mathbf{P}_{0}\right|_{\sigma(X_{0},X_{1},\dots, X_{n})}=  %\frac{\nu_{X_0}}{\nu_{0,X_0}}   
    \exp\left(\sum_{k=1}^{n} \log \frac{P_{k}(X_{k-1},X_{k})}{P_{0}(X_{k-1},X_{k})}
    \right). \tag{D}    \end{equation*}
    Conversely, if there exists $\mathbf{P}_{0}$ as described above and $\mathbf{P}\ll \mathbf{P}_{0}$, then there exists a $\sigma(X_{0},X_{1},\dots,X_{n})$-predictable transition probability matrix $P_n$ such that $X_n$ is a Markov chain with transition probability matrix $P_n$ under $\mathbf P$ and (D) holds. 
\end{theorem}

\begin{proof} Sufficiency. By uniqueness of the martingale problem for discrete-time Markov chains, it suffices to prove $\mathbf P \ll \mathbf P_0$ if $X_n$ has the Markov property under both $\mathbf P$ and $\mathbf P_0$, and is stationary under $\mathbf P_0$.  Let 
\[Z_n = \exp\left(\sum_{k=1}^{n} \log \frac{P_{k}(X_{k-1},X_{k})}{P_{0}(X_{k-1},X_{k})}
    \right),\]
Then $Z_n$ is a $\mathbf P_0$-martingale by the stationarity of $X_n$ under $\mathbf P_0$, which can be seen via the followiong calculation
$$\mathbf E_{0}\left[\frac{P_{n}(X_{n-1},X_{n})}{P_{0}(X_{n-1},X_{n})}\Bigm|\mathscr F_{n-1}\right]   =\sum_{i=1}^N\frac{P_{n}(X_{n-1},i)}{P_{0}(X_{n-1},i)}P_0(X_{n-1},i)=1.   $$
Define a probability measure $\mathbf P'$ on $\sigma(X_0,\dots,X_n)$ by $d\mathbf P' = Z_n \, d\mathbf P_0$. Intuitively, $\mathbf P$ and $\mathbf P'$ coincide on $\sigma(X_{0},\dots,X_{n})$. Since $Z_n$ is bounded and hence uniformly integrable, this identification extends to the full $\sigma$-field $\sigma\big(\bigcup_n \sigma(X_{0},\dots,X_{n})\big)$. Thus $\mathbf P = \mathbf P'$, and $(D)$ holds.

%Necessity. We need the following lemma.
    \begin{lemma}
        Let $(X_n, P_n)$ be an adapted process, where $P_n$ is an $\mathscr F_{n-1}$-measurable random transition probability matrix with row sum 1. In order that for every $z \in \mathbb R^N$ the process
        \begin{equation}
            M_n^z = z_{X_n} - z_{X_0} - \sum_{k=1}^{n} \big( (P_k - I)z\big)_{X_{k-1}}
            \label{eq: Markov martingale}
        \end{equation}
        is an  $\mathscr F_n$-martingale, it is sufficient and necessary that $X_n$ is a Markov chain with transition matrix $P_n$ with respect to $\mathscr F_n$. 
        \end{lemma}
    \begin{proof} Suppose $X_n$ is a Markov chain with transition matrices $P_n$ with respect to $\mathscr F_n$. Then for every $n$ and $j$,
    $$\mathbf P(X_{n} = j \mid \mathscr F_{n-1}) = P_{n, X_{n-1}, j}.$$
    Or equivalently, for every $z \in \mathbb R^N$,
    $$\mathbf E[z_{X_{n}} \mid \mathscr F_{n-1}] = (P_{n} z)_{X_{n-1}}.$$
    Now consider $M_n^z$. To show it is a martingale, we compute
    $$M_{n}^z = z_{X_{n}} - z_{X_0} - \sum_{k=1}^{n} \big( (P_k-I) z\big)_{X_{k-1}}.$$
    Thus,
    $$M_{n}^z - M_{n-1}^z = z_{X_{n}} - z_{X_{n-1}} -\big( (P_{n}-I) z\big)_{X_{n-1}} = z_{X_{n}} - (P_{n} z)_{X_{n-1}}.$$
    Taking conditional expectation,
    $$\mathbf E[M_{n}^z - M_{n-1}^z \mid \mathscr F_{n-1}] = \mathbf E[z_{X_{n}} \mid \mathscr F_{n-1}] - (P_{n} z)_{X_{n-1}} = 0.$$
    Hence, $\mathbf E[M_{n}^z \mid \mathscr F_{n-1}] = M_{n-1}^z$, so $M_n^z$ is a martingale.

    Conversely, suppose $M_n^z$ is a martingale for every $z \in \mathbb R^N$. Then
    $$\mathbf E[M_{n}^z - M_{n-1}^z \mid \mathscr F_{n-1}] = 0.$$
    But
    $$M_{n}^z - M_{n-1}^z = z_{X_{n}} - (P_{n} z)_{X_{n-1}},$$
    so
    $$\mathbf E[z_{X_{n}} \mid \mathscr F_{n-1}] = (P_{n} z)_{X_{n-1}}.$$
    This holds for all $z \in \mathbb R^N$. In particular, choose $z = (z_1,z_2,\dots,z_N)$ to be the indicator vector: $z_i = \mathbf 1{\{i=j\}}$. Then $z_{X_{n}} = \mathbf 1{\{X_{n} = j\}}$, and
    $$(P_{n} z)_{X_{n-1}} = P_{n, X_{n-1}, j}.$$
    Thus,
    $$\mathbf P(X_{n} = j \mid \mathscr F_{n-1}) = P_{n, X_{n-1}, j}.$$
    This is exactly the Markov property with transition matrices $P_n$ with respect to $\mathscr F_n$. \end{proof}


    Necessity, Since $\mathbf{P} \ll \mathbf{P}_0$, we define the likelihood ratio process
    $$Z_n = \mathbf E_0 \left[ \frac{d\mathbf{P}}{d\mathbf{P}_0} \Bigm| {\sigma(X_0, X_1, \dots, X_n)}\right].$$
    The process $Z_n$ is a $\mathbf{P}_0$-martingale, and $Z_n > 0$ $\mathbf{P}_0$-a.s. For a fixed \( n \ge 1 \), define for each \( j \),
    \[    \Delta_{n,j} := \mathbf{1}{\{X_n = j\}} - P_0(X_{n-1}, j).\]

    \begin{lemma}\label{lem: delta-basis} For every \( \mathscr{F}_n \)-measurable random variable \( Y_n \) satisfying \( \mathbf{E}_0[Y_n \mid \mathscr{F}_{n-1}] = 0 \), there exist \( \mathscr{F}_{n-1} \)-measurable coefficients \( G_{n,j} \) such that
    \[Y_n = \sum_{j=1}^N G_{n,j}\Delta_{n,j}, \qquad \mathbf{P}_0\text{-a.s.}.\]
    and on each atom \( \{X_{n-1} = i\} \), the coefficients \( G_{n,j} \) are uniquely determined by a linear system.
    \end{lemma}

\begin{proof}
Consider an atom \( A_i := \{X_{n-1} = i\} \) for some \( i\). On \( A_i \), define the random vector \( (Y_n(i,j))_{j=1}^N \) by setting \( Y_n(i,j) = Y_n \cdot \mathbf{1}{\{X_n = j\}} \) for each \( j \). Note that \( Y_n(i,j) \) represents the value of \( Y_n \) on the event \( \{X_n = j\} \cap A_i \). The condition \( \mathbf{E}_0[Y_n \mid \mathscr{F}_{n-1}] = 0 \) implies that
\[\sum_{j=1}^N  P_0(i,j) Y_n(i,j) = 0,\]
since the conditional expectation given \( \mathscr{F}_{n-1} \) reduces to averaging over \( X_n \) with weights \( P_0(i,j) \) on \( A_i \).

Now, consider the vectors \( \Delta_{n,j}(i, \cdot) \) for \(j\), where \( \Delta_{n,j}(i, k) = \mathbf{1}{\{k = j\}} - P_0(i,j) \) for each \(k\). These vectors span the subspace
\[V_i = \left\{ v \in \mathbb{R}^N : \sum_{j=1}^N  P_0(i,j) v_j = 0 \right\},\]
because any vector \( v \in V_i \) can be written as a linear combination of the vectors \( \Delta_{n,j}(i, \cdot) \). Specifically, the system of equations
\[Y_n(i, k) = \sum_{j=1}^N  G_{n,j}(i) \left[ \mathbf{1}{\{k = j\}} - P_0(i,j) \right], \qquad k = 1,2,\dots, N,\]
has a solution for coefficients \( G_{n,j}(i) \) since the right-hand side covers all vectors in \( V_i \). The solution is unique up to the constraints of the subspace.

Since this holds for each atom \( A_i \), we can define \( \mathscr{F}_{n-1} \)-measurable coefficients \( G_{n,j} \) by setting \( G_{n,j} = G_{n,j}(i) \) on \( A_i \). Then, the decomposition
\[Y_n = \sum_{j=1}^N G_{n,j} \, \Delta_{n,j}\]
holds \( \mathbf{P}_0 \)-almost surely, as required.
\end{proof}

For each $k \leq n$, set $Y_k = Z_k - Z_{k-1}$. Since $(Z_k)$ is a martingale, $\mathbf{E}_0[Y_k \mid \mathscr{F}_{k-1}] = 0$. By Lemma \ref{lem: delta-basis}, there exist $\mathscr{F}_{k-1}$-measurable coefficients $(G_{k,j})_{j=1}^N$ such that
$$Y_k = \sum_{j=1}^N G_{k,j}\,\Delta_{k,j}, 
\qquad \Delta_{k,j} = \mathbf{1}{\{X_k=j\}} - P_0(X_{k-1},j),$$
which holds $\mathbf P_0$-almost surely. Intuitively, the variables $\Delta_{k,j}$ represent centered indicators of the transition into state $j$, so the above identity is a martingale representation of $Y_k$. 

Next, we identify the transition probabilities $P_k(i,j)$ under $\mathbf P$. By the likelihood ratio property, we have under $\mathbf P_0$ the recursion
$$Z_k = Z_{k-1}\frac{P_{k}(X_{k-1},X_k)}{P_0(X_{k-1},X_k)},$$
so that comparing this expression with the martingale representation of $Y_k$ and restricting to the event $\{X_{k-1}=i,\,X_k=j\}$,
$$Z_{k-1}\!\left(\frac{P_{k}(i,j)}{P_0(i,j)}-1\right) = G_{k,j} - \sum_{l=1}^N G_{k,l}P_0(i,l).$$
Therefore, whenever $Z_{k-1}>0$, the transition probabilities under $\mathbf P$ are given explicitly by
$$P_{k}(i,j) = P_0(i,j)\left(1+\frac{1}{Z_{k-1}}\Big(G_{k,j}-\sum_{l=1}^N G_{k,l}P_0(i,l)\Big)\right).$$

It remains to check that these coefficients indeed form a stochastic matrix. Fixing $i$ and summing over $j$, we obtain
$$
\sum_{j=1}^N P_{k}(i,j)
= \sum_{j=1}^N P_0(i,j) 
+ \frac{1}{Z_{k-1}}\!\left(\sum_{j=1}^N P_0(i,j)G_{k,j} - \sum_{l=1}^N G_{k,l}P_0(i,l)\right).
$$
The second term vanishes because both expressions inside coincide, leaving
$$\sum_{j=1}^N P_0(i,j)  = 1.$$
Thus each row sums to one, and non-negativity holds by construction, since the $P_k(i,j)$ arise as conditional probabilities. Consequently, $P_k$ is a valid transition matrix adapted to the natural filtration. In particular, the process $X_n$ is a Markov chain under $\mathbf{P}$, and the Girsanov transform (D) follows.  This complete the proof.  \end{proof}

\end{document}
